\section{Mathematical notation and equivariant operators}
\label{sec:sm-model-formulation}
\label{sec:sm-methods}

This section collects the mathematical details that supplement the architecture description in Section~\ref{sec:methods}: the formal SO(3) representation notation, the edge-local frame and SO(2) decomposition, the closed-form cutoff envelope coefficients, the truncated SO(2) layout and rescaled lift, the SO(2)-equivariant operator algebra, and the cutoff-consistent first-layer bias correction used by the ablations in later supplementary sections.

\subsection{SO(3) representation: notation}
\label{sec:sm-representation}

For each integer $l\ge 0$, let $V_l\cong\mathbb{R}^{2l+1}$ be the real $(2l+1)$-dimensional irreducible representation of SO(3), spanned by the real spherical harmonics $\{Y_l^m\}_{m=-l}^{l}$.
The action of a rotation
$Q\in\mathrm{SO}(3)$ on $V_l$ is given by the real Wigner $D$-matrix
$D^l(Q)\in\mathrm{O}(2l+1)$, the orthogonal $(2l+1)\times(2l+1)$ matrix
that describes how the real spherical harmonics transform under rotation,
\begin{equation}
    Y_l^m(Q^{-1}\widehat{\mathbf{r}}) =\sum_{m'=-l}^{l}D^l(Q)_{m,m'}\,Y_l^{m'}(\widehat{\mathbf{r}}).
    \label{eq:sm-wigner-D}
\end{equation}
The map $D^l:\mathrm{SO}(3)\to\mathrm{O}(2l+1)$ is a continuous group homomorphism (i.e.\ $D^l(QQ')=D^l(Q)D^l(Q')$), and the $V_l$ are \emph{irreducible}: $V_l$ admits no proper SO(3)-invariant subspace.
The trivial case $l=0$ is the invariant scalar representation with $D^0(Q)\equiv 1$.

The node-feature space $V_{\le L}\otimes\mathbb{R}^{C}=\bigoplus_{l=0}^{L}
    V_l\otimes\mathbb{R}^{C}$ of main-text Section~\ref{sec:methods} therefore
carries the real block-diagonal action
\begin{equation}
    D(Q)=\bigoplus_{l=0}^{L}D^l(Q), \label{eq:sm-direct-sum} \end{equation} acting independently on each degree-$l$ block and trivially on the channel factor $\mathbb{R}^{C}$.
Basis coefficients are packed by the linear index
\begin{equation}
    \iota(l,m)=l^2+l+m,\qquad
    m=-l,\ldots,l,\quad l=0,\ldots,L,
    \label{eq:sm-pack-index}
\end{equation}
used throughout this Supplement and in main-text Eq.~\eqref{eq:dpa4-gie};
under this packing, the matrix $D(Q)$ is block-diagonal with the
$(2l+1)\times(2l+1)$ block $D^l(Q)$ occupying rows and columns
$\iota(l,-l),\ldots,\iota(l,+l)$.
A map $f:V_{\le L}\otimes\mathbb{R}^{C}\to V_{\le L}\otimes\mathbb{R}^{C'}$ is \emph{SO(3)-equivariant} if $f\circ D(Q)=D(Q)\circ f$ for every $Q\in\mathrm{SO}(3)$, and \emph{SO(3)-invariant} if it lands in the $l=0$ block, on which $D^0\equiv 1$.

\subsection{Edge-local frame and SO(2) decomposition}
\label{sec:sm-local-frame}

\paragraph{Goal.}
For each directed edge $(i,j)$, DPA4 chooses a rotation $R_{ij}$ satisfying
\begin{equation}
    R_{ij}\widehat{\mathbf{r}}_{ij}=\mathbf{e}_z=(0,0,1)^{\top}.
    \label{eq:sm-local-frame-goal}
\end{equation}
This converts an SO(3) problem into an SO(2) problem in the edge-local frame: after the bond direction is aligned with $\mathbf{e}_z$, the residual symmetry is the subgroup of rotations about $\mathbf{e}_z$.
The in-plane basis is a gauge choice.
The equivariant message does not depend on this choice because the local operator commutes with the residual SO(2) action.

\paragraph{Two quaternion charts.}
Let $\widehat{\mathbf{r}}=(x,y,z)\in \Sph$.
DPA4 uses two smooth unit-quaternion
charts,
\begin{equation}
    \mathbf{q}^{+}(\widehat{\mathbf{r}})
    =
    \frac{(1+z,\;y,\;-x,\;0)}{\sqrt{2(1+z)}},
    \qquad
    \mathbf{q}^{-}(\widehat{\mathbf{r}})
    =
    \frac{(-x,\;0,\;1-z,\;y)}{\sqrt{2(1-z)}}.
    \label{eq:sm-quaternion-charts}
\end{equation}
The first chart is regular away from the south pole and the second is regular away from the north pole.
In the overlap, the sign of $\mathbf{q}^{-}$ is chosen
so that $\langle\mathbf{q}^{+},\mathbf{q}^{-}\rangle\ge 0$, and the blended
quaternion is
\begin{equation}
    \mathbf{q}_{ij}
    =
    \frac{\lambda\mathbf{q}^{+}(\widehat{\mathbf{r}}_{ij})
        +(1-\lambda)\mathbf{q}^{-}(\widehat{\mathbf{r}}_{ij})}
    {\left\|\lambda\mathbf{q}^{+}(\widehat{\mathbf{r}}_{ij})
        +(1-\lambda)\mathbf{q}^{-}(\widehat{\mathbf{r}}_{ij})\right\|},
    \qquad
    \lambda=\frac{1+z}{2}.
    \label{eq:sm-quaternion-blend}
\end{equation}
The denominator is bounded away from zero after shortest-arc sign alignment in the chosen chart overlap, so the resulting gauge is smooth on the chart used by the implementation.
The Wigner-D matrix in the edge-local gauge is denoted $D_{ij}=D(R(\mathbf{q}_{ij}))$.
The underlying unit-quaternion rotation matrix is
\begin{equation}
    R(\mathbf{q})=
    \begin{pmatrix}
        1-2(q_y^2+q_z^2) & 2(q_xq_y-q_wq_z) & 2(q_xq_z+q_wq_y) \\
        2(q_xq_y+q_wq_z) & 1-2(q_x^2+q_z^2) & 2(q_yq_z-q_wq_x) \\
        2(q_xq_z-q_wq_y) & 2(q_yq_z+q_wq_x) & 1-2(q_x^2+q_y^2)
    \end{pmatrix}.
    \label{eq:sm-quaternion-rotation}
\end{equation}
During training, an additional random roll about the local $z$ axis may be composed with $R_{ij}$.
This roll is a gauge augmentation: it changes the in-plane basis of the local gauge but leaves the bond direction fixed.
Because the SO(2) stack is gauge equivariant, the lifted global message is unchanged by this roll apart from the prescribed SO(3) transformation law.

\paragraph{SO(2) decomposition.}
Let $\{\mathbf{e}^{(l)}_m\}_{m=-l}^{l}$ denote the real-spherical-harmonic basis of $V_l$.
Under restriction to the SO(2) subgroup of rotations about
the local $z$ axis, $V_l$ decomposes as
\begin{equation}
    V_l|_{\mathrm{SO}(2)} = \mathrm{span}_{\mathbb{R}}\{\mathbf{e}^{(l)}_0\} \;\oplus\; \bigoplus_{m=1}^{l} \mathrm{span}_{\mathbb{R}}\{\mathbf{e}^{(l)}_{-m},\,\mathbf{e}^{(l)}_{+m}\}, \label{eq:sm-so2-decomposition} \end{equation} where $\mathrm{span}_{\mathbb{R}}\{\cdot\}$ denotes the real linear span of the indicated basis vectors.
The $m=0$ summand is a one-dimensional trivial SO(2) sub-representation, and each $m=1,\ldots,l$ summand is a two-dimensional real SO(2) sub-representation on which a $z$-axis rotation by angle $\theta$ acts as planar rotation by angle $m\theta$, equivalent to multiplication by the complex phase $e^{im\theta}$.
This is the algebraic reason DPA4 can use SO(2)-equivariant local operators instead of Clebsch--Gordan tensor products.

\subsection{Closed-form cutoff coefficients}
\label{sec:sm-radial}

The four boundary conditions
$s_p(r_{\mathrm{c}})=s_p'(r_{\mathrm{c}})=s_p''(r_{\mathrm{c}})=s_p'''(r_{\mathrm{c}})=0$
of the cutoff envelope $s_p(r)$ in main-text
Eq.~\eqref{eq:dpa4-envelope} uniquely fix the coefficients to
\begin{align}
    a_p & =-\frac{(p+1)(p+2)(p+3)}{6}, &
    b_p & =\frac{p(p+2)(p+3)}{2},        \\
    c_p & =-\frac{p(p+1)(p+3)}{2},     &
    d_p & =\frac{p(p+1)(p+2)}{6}.
    \label{eq:sm-cutoff-coefficients}
\end{align}
For the two values used in DPA4 ($s_5$ for edge weighting and the smooth
degree, $s_7$ inside the radial basis of main-text
Eq.~\eqref{eq:dpa4-radial-basis}), the explicit polynomials are
\begin{align}
    s_5(r) & =1-56x^5+140x^6-120x^7+35x^8,       \\
    s_7(r) & =1-120x^7+540x^8-1080x^9+840x^{10},
\end{align}
with $x=r/r_{\mathrm{c}}$.

\subsection{Truncated local layout and rescaled lift}
\label{sec:sm-truncation}

Inside the edge-local frame, DPA4 retains only
\begin{equation}
    \mathcal{I}_M=
    \{(l,m):0\le l\le L,\ |m|\le \min(l,M)\},
    \qquad
    D_M=|\mathcal{I}_M|.
    \label{eq:sm-trunc-index}
\end{equation}
Let $P_M:\mathbb{R}^{(L+1)^2}\to\mathbb{R}^{D_M}$ be the orthogonal projection onto this reduced layout.
The truncated edge rotation is
\begin{equation}
    D_{ij}^{\le M}=P_M D_{ij}.
    \label{eq:sm-truncated-rotation}
\end{equation}
Because $P_M^{\top}P_M$ is not the identity when $M<L$, the round trip loses degree-block norm.
DPA4 applies the diagonal lift compensation
\begin{equation}
    (\Xi_M)_{\iota(l,m),\iota(l,m)}
    =
    \kappa_l
    =
    \sqrt{\frac{2l+1}{2\min(l,M)+1}}.
    \label{eq:sm-rescale}
\end{equation}
Thus $\Xi_M= \operatorname{diag}(\kappa_0 I_1,\kappa_1 I_3,\ldots,\kappa_L I_{2L+1})$.

\subsection{Equivariant operator algebra}
\label{sec:sm-operators}

\paragraph{Classification by Schur's lemma.}
The linear operators used by DPA4 follow from the classification of equivariant maps on the relevant representation spaces.
In the global
SO(3) frame, Schur's lemma forbids mixing distinct degrees:
\begin{equation}
    \mathrm{Hom}_{\mathrm{SO}(3)}
    (V_{\le L}\otimes\mathbb{R}^{C},V_{\le L}\otimes\mathbb{R}^{C'})
    \cong
    \bigoplus_{l=0}^{L}\mathbb{R}^{C\times C'}.
    \label{eq:sm-schur-so3}
\end{equation}
In the edge-local SO(2) frame, the $m=0$ lines are trivial representations
and each $|m|>0$ pair is complex type, giving
\begin{equation}
    \mathrm{Hom}_{\mathrm{SO}(2)}
    (V_{\le L}\otimes\mathbb{R}^{C},V_{\le L}\otimes\mathbb{R}^{C'})
    \cong
    \bigoplus_{l,l'}\mathbb{R}^{C\times C'}
    \oplus
    \bigoplus_{m=1}^{M}\bigoplus_{l,l'\ge m}\mathbb{C}^{C\times C'}.
    \label{eq:sm-schur-so2}
\end{equation}
Equation~\eqref{eq:sm-schur-so2} is the formal reason the local-frame operator may mix degrees $l,l'$ while preserving each $|m|$ stratum.

\paragraph{Global SO(3) linear maps.}
By Eq.~\eqref{eq:sm-schur-so3}, an SO(3)-equivariant \emph{degree-wise} map
has the explicit form
\begin{equation}
    (L^{\mathrm{deg}}_{\Theta}\mathbf{h})_{\iota(l,m),c'}
    =
    \sum_{c=1}^{C}
    W^{(l)}_{c,c'}\mathbf{h}_{\iota(l,m),c}, \label{eq:sm-degree-linear} \end{equation} with one learnable channel matrix $W^{(l)}\in\mathbb{R}^{C\times C'}$ per degree $l$.
Here ``degree-wise'' means \emph{per-$l$ but identical across all $m\in\{-l,\ldots,l\}$ within a degree}: the same $W^{(l)}$ is applied to every $m$-slice of the degree-$l$ block, with different matrices allowed for different $l$.
All ``degree-wise channel-mixing maps'' in main-text Section~\ref{sec:methods} (including $L^{\mathrm{pre}}_{\mathrm{deg}}$, $L^{\mathrm{post}}_{\mathrm{deg}}$, $L^{\mathrm{rad}}_{\mathrm{lift}}$, $L^{\mathrm{ch}}_{\mathrm{in}}$ and $L^{\mathrm{ch}}_{\mathrm{out}}$) are instances of this form.

\paragraph{Edge-local SO(2) linear maps.}
In the edge-local frame, the $m=0$ lines are trivial SO(2) representations, whereas each $|m|>0$ pair is complex type.
The most general SO(2)-equivariant
linear map on the reduced layout is therefore
\begin{align}
    (L_{\Theta}^{\mathrm{SO2}}\mathbf{x})_{(l,0),c'}
     & =
    \sum_{l'=0}^{L}\sum_{c=1}^{C}
    A_{c,c'}^{(l,l',0)}\mathbf{x}_{(l',0),c}
    +b_{0,c'}\delta_{l,0},
    \label{eq:sm-so2-m0} \\
    \begin{pmatrix}
        (L_{\Theta}^{\mathrm{SO2}}\mathbf{x})_{(l,-m),c'}\\
        (L_{\Theta}^{\mathrm{SO2}}\mathbf{x})_{(l,+m),c'}
    \end{pmatrix}
     & =
    \sum_{l'\ge m}\sum_{c=1}^{C}
    \begin{pmatrix}
        U_{c,c'}^{(l,l',m)} & -V_{c,c'}^{(l,l',m)} \\
        V_{c,c'}^{(l,l',m)} & U_{c,c'}^{(l,l',m)}
    \end{pmatrix}
    \begin{pmatrix}
        \mathbf{x}_{(l',-m),c} \\
        \mathbf{x}_{(l',+m),c}
    \end{pmatrix}.
    \label{eq:sm-so2-mn}
\end{align}

\paragraph{Equivariant RMS normalization.}
For $\mathbf{h}\in V_{\le L}\otimes\mathbb{R}^{C}$, define
\begin{equation}
    \sigma^2(\mathbf{h})
    =
    \sum_{l=0}^{L}\sum_{m=-l}^{l}\sum_{c=1}^{C}
    \frac{(\mathbf{h}_{\iota(l,m),c}-\delta_{l,0}\bar{h})^2}
    {(2l+1)(L+1)C},
    \qquad
    \bar{h}=C^{-1}\sum_c \mathbf{h}_{\iota(0,0),c}.
    \label{eq:sm-rms-stat}
\end{equation}
The normalization
\begin{equation}
    (N_{\gamma,\beta}\mathbf{h})_{\iota(l,m),c}
    =
    \gamma_l
    \frac{\mathbf{h}_{\iota(l,m),c}-\delta_{l,0}\bar{h}}
    {\sqrt{\sigma^2(\mathbf{h})+\varepsilon}}
    +\delta_{l,0}\beta_c
    \label{eq:sm-rmsnorm}
\end{equation}
commutes with SO(3), because the numerator is equivariant, the denominator is
invariant under the orthogonal representation $D^l$, and $\gamma_l$ is constant
within each degree block.

\paragraph{Scalar-gated nonlinearity.}
For a smooth scalar nonlinearity $\psi$ and degree-wise gate matrices
$G^{(l)}$, define
\begin{equation}
    (\Gamma_{\psi,G}\mathbf{h})_{\iota(l,m),c}
    =
    \begin{cases}
        \psi(\mathbf{h}_{\iota(0,0),c}),                  & l=0,    \\[2pt]
        \mathbf{h}_{\iota(l,m),c}
        \sigma\!\left(\sum_{c'}
        G^{(l)}_{c',c} \mathbf{h}_{\iota(0,0),c'}\right), & l\ge 1.
    \end{cases}
    \label{eq:sm-gated}
\end{equation}
The gate is a scalar function of the invariant $l=0$ slice, so the operation is SO(3)-equivariant.

\paragraph{Scalar-keyed mixtures.}
If $\mathbf{x}^{(1)},\ldots,\mathbf{x}^{(S)}$ are equivariant tensors of the
same type and $\pi$ is an invariant scalar projection, then
\begin{equation}
    \mathcal{A}(\mathbf{x}^{(1)},\ldots,\mathbf{x}^{(S)};\mathbf{q})
    =
    \sum_{s=1}^{S}\alpha_s\mathbf{x}^{(s)},\qquad
    \alpha_s=
    \frac{\exp\langle\mathbf{q},\pi(\mathbf{x}^{(s)})\rangle}
    {\sum_{s'}\exp\langle\mathbf{q},\pi(\mathbf{x}^{(s')})\rangle}
    \label{eq:sm-scalar-keyed-mixture}
\end{equation}
is equivariant, because the weights are invariant scalars.

\subsection{Bias consistency at the smooth cutoff}
\label{sec:sm-so2conv}

If the first SO(2) linear map of the EMFA SO(2) convolution (main-text Sec.~\ref{sec:so2}; Eqs.~\eqref{eq:sm-so2-m0}, \eqref{eq:sm-so2-mn}) includes an additive $l=0$ bias, that bias must vanish with the same smooth envelope as the rest of the edge message.
Otherwise an edge whose radial contribution has gone to zero could still carry a constant scalar offset.
DPA4 therefore treats the first-layer scalar bias as part of the edge-conditioned response.
If $b_{0,f,c}$ is the scalar bias for focus $f$
and channel $c$, the net contribution to the local $(l,m)=(0,0)$ slot is
adjusted to
\begin{equation}
    b_{0,f,c}\,\widetilde{\rho}_{ij,0,c}\,s_5(r_{ij}),
    \label{eq:sm-bias-consistency}
\end{equation}
which is $C^3$ and vanishes at the cutoff.
This correction is only needed at the first SO(2) layer because subsequent layers operate on already modulated local features.

\subsection{Numerical equivariance of truncated local layouts}
\label{sec:sm-truncated-equivariance}

Table~\ref{tab:sm_s2_truncated_equivariance} extends the full-coefficient comparison of main-text Table~\ref{tab:s2_full_equivariance} to the $m$-truncated local layout used inside the EMFA SO(2) convolution, evaluating the maximum equivariance error of product-grid rules and Lebedev quadrature under random rotations about the local $z$ axis.

\begin{table}[p]
    \centering
    \scriptsize
    \setlength{\tabcolsep}{3pt}
    \caption{
        $m$-truncated \texorpdfstring{$\Sph$}{S2} activation equivariance under random local
        $z$-axis rotations.
        These cases test the reduced local layout used in SO(2) convolution.
        Product-grid rules are $(R_{\phi},R_{\theta})$ with the total number of grid points $R_{\phi}R_{\theta}$ given alongside; Lebedev rules are reported by their algebraic order of accuracy $p$ and the corresponding number of points.
        Errors are maximum absolute deviations between the two equivariance paths.
    }
    \label{tab:sm_s2_truncated_equivariance}
    \begin{tabular*}{\linewidth}{@{\extracolsep{\fill}}cccccccccc@{}}
        \toprule
        \textbf{$M$} &
        \textbf{$L$} &
        \multicolumn{4}{c}{\textbf{Product grid}} &
        \multicolumn{4}{c}{\textbf{Lebedev quadrature}} \\
        \cmidrule(lr){3-6}\cmidrule(lr){7-10}
        & & \textbf{Rule} & \textbf{\#\,pts} & \textbf{fp64 error} & \textbf{fp32 error}
        & \textbf{$p$} & \textbf{\#\,pts} & \textbf{fp64 error} & \textbf{fp32 error} \\
        \midrule
        1 & 2 & $6\times8$   & 48  & $2.36\times10^{-7}$ & $3.58\times10^{-7}$ & 7  & 26  & $2.31\times10^{-14}$ & $2.38\times10^{-7}$ \\
        1 & 3 & $6\times12$  & 72  & $1.22\times10^{-7}$ & $5.96\times10^{-7}$ & 9  & 38  & $3.55\times10^{-14}$ & $2.98\times10^{-7}$ \\
        1 & 4 & $6\times14$  & 84  & $1.12\times10^{-6}$ & $9.54\times10^{-7}$ & 13 & 74  & $1.04\times10^{-13}$ & $9.54\times10^{-7}$ \\
        1 & 5 & $6\times18$  & 108 & $1.10\times10^{-7}$ & $1.43\times10^{-6}$ & 15 & 86  & $9.34\times10^{-14}$ & $7.15\times10^{-7}$ \\
        1 & 6 & $6\times20$  & 120 & $7.64\times10^{-7}$ & $1.91\times10^{-6}$ & 19 & 146 & $8.56\times10^{-14}$ & $2.15\times10^{-6}$ \\
        1 & 7 & $6\times24$  & 144 & $2.17\times10^{-7}$ & $1.91\times10^{-6}$ & 21 & 170 & $2.08\times10^{-13}$ & $3.34\times10^{-6}$ \\
        \midrule
        2 & 2 & $8\times8$   & 64  & $4.01\times10^{-7}$ & $8.34\times10^{-7}$ & 7  & 26  & $1.50\times10^{-14}$ & $2.38\times10^{-7}$ \\
        2 & 3 & $8\times12$  & 96  & $5.99\times10^{-7}$ & $8.34\times10^{-7}$ & 9  & 38  & $5.71\times10^{-14}$ & $3.58\times10^{-7}$ \\
        2 & 4 & $8\times14$  & 112 & $6.02\times10^{-7}$ & $1.67\times10^{-6}$ & 13 & 74  & $9.15\times10^{-14}$ & $5.96\times10^{-7}$ \\
        2 & 5 & $8\times18$  & 144 & $1.19\times10^{-6}$ & $1.55\times10^{-6}$ & 15 & 86  & $7.83\times10^{-14}$ & $4.77\times10^{-7}$ \\
        2 & 6 & $8\times20$  & 160 & $1.33\times10^{-6}$ & $2.15\times10^{-6}$ & 19 & 146 & $1.29\times10^{-13}$ & $9.54\times10^{-7}$ \\
        2 & 7 & $8\times24$  & 192 & $1.41\times10^{-6}$ & $2.62\times10^{-6}$ & 21 & 170 & $1.56\times10^{-13}$ & $1.43\times10^{-6}$ \\
        \bottomrule
    \end{tabular*}
\end{table}

\subsection{Completeness of the Wigner bilinear network}
\label{sec:sm-wbn-complete}

This subsection proves the completeness statement invoked in main-text Section~\ref{sec:s2-method}: networks stacking Wigner bilinear blocks with the \emph{full} frame range $|k|\le l$, interleaved with equivariant linear maps, approximate every continuous SO(3)-equivariant map between band-limited feature spaces, and realize polynomial equivariant maps exactly.
The block analyzed here is the exact continuous model underlying the deployed FFN: it retains all frames, forms products over all pairs of left and right channels, and evaluates the Haar integral exactly.
The deployed block (main-text Eqs.~\eqref{eq:dpa4-wbb-lift}--\eqref{eq:dpa4-ffn}) is its economical specialization -- channel-diagonal products, frames truncated to $|k|\le k_{\max}$, and an exact quadrature in place of the integral -- and the role of the truncation is discussed in Remark~\ref{rem:sm-wbn-truncation}.

Throughout, $g$ denotes a rotation, $dg$ the normalized Haar measure on $SO(3)$, and $D^{(l)}_{mk}(g)$ the \emph{complex} Wigner matrices, which satisfy the orthogonality relation
\begin{equation}
    (2l+1)\int_{SO(3)}\overline{D^{(l)}_{mk}(g)}\,D^{(l')}_{m'k'}(g)\,dg
    =\delta_{ll'}\delta_{mm'}\delta_{kk'}
    \label{eq:sm-wigner-orth}
\end{equation}
and the product identity~\cite{varshalovich1988quantum}
\begin{equation}
    D^{(l_1)}_{m_1k_1}(g)\,D^{(l_2)}_{m_2k_2}(g)
    =\sum_{L=|l_1-l_2|}^{l_1+l_2}
    \langle l_1 m_1\, l_2 m_2|L M\rangle\,
    \langle l_1 k_1\, l_2 k_2|L \kappa\rangle\,
    D^{(L)}_{M\kappa}(g),
    \label{eq:sm-wigner-product}
\end{equation}
with $M=m_1+m_2$, $\kappa=k_1+k_2$, and $\langle\cdot\,\cdot|\cdot\rangle$ the Clebsch--Gordan coefficients.
The associated coupling map $C^{L}_{l_1l_2}:V_{l_1}\otimes V_{l_2}\to V_L$,
\begin{equation}
    \bigl(C^{L}_{l_1l_2}(x,y)\bigr)_M
    =\sum_{m_1,m_2}\langle l_1 m_1\, l_2 m_2|L M\rangle\,x_{m_1}y_{m_2},
    \label{eq:sm-cg-map}
\end{equation}
is a nonzero SO(3)-intertwiner for every admissible triple $|l_1-l_2|\le L\le l_1+l_2$.
The deployed network uses the real Wigner basis of Section~\ref{sec:sm-representation}; the real and complex conventions are related by an invertible equivariant change of coordinates, which is absorbed into the learnable channel maps, so it suffices to argue in the complex convention.

\paragraph{The full-frame block.}
Let $\mathcal{H}=V_{\le L_{\max}}\otimes\mathbb{C}^{C}$.
A \emph{full-frame Wigner bilinear block} $\mathcal{B}:\mathcal{H}\to\mathcal{H}$ maps a feature $x^{(l)}_{m,c}$ to
\begin{align}
    u^{\mathrm{L},(l)}_{m,k,a}
    &=\sum_{c}A^{\mathrm{L},(l)}_{k,a,c}\,x^{(l)}_{m,c},
    \qquad
    u^{\mathrm{R},(l)}_{m,k,b}
    =\sum_{c}A^{\mathrm{R},(l)}_{k,b,c}\,x^{(l)}_{m,c},
    \qquad |k|\le l,
    \label{eq:sm-wbb-lin}\\
    f^{\mathrm{L}}_{a}(g)
    &=\sum_{l,m,k}u^{\mathrm{L},(l)}_{m,k,a}D^{(l)}_{mk}(g),
    \qquad
    f^{\mathrm{R}}_{b}(g)
    =\sum_{l,m,k}u^{\mathrm{R},(l)}_{m,k,b}D^{(l)}_{mk}(g),
    \label{eq:sm-wbb-synth}\\
    h_{ab}(g)
    &=f^{\mathrm{L}}_{a}(g)\,f^{\mathrm{R}}_{b}(g),
    \label{eq:sm-wbb-prod}\\
    z^{(l)}_{m,k,ab}
    &=(2l+1)\int_{SO(3)}\overline{D^{(l)}_{mk}(g)}\,h_{ab}(g)\,dg,
    \label{eq:sm-wbb-analysis}\\
    \mathcal{B}(x)^{(l)}_{m,d}
    &=\sum_{c}P^{(l)}_{dc}\,x^{(l)}_{m,c}
    +\delta_{l0}\delta_{m0}\,q_{d}
    +\sum_{k,a,b}B^{(l)}_{d,k,ab}\,z^{(l)}_{m,k,ab},
    \label{eq:sm-wbb-out}
\end{align}
with the per-degree matrices $A^{\mathrm{L}},A^{\mathrm{R}},P$, the bias $q$, and the mixing coefficients $B$ learnable.
A \emph{Wigner bilinear network} (WBN) is a composition $\pi\circ\mathcal{B}_{K}\circ\cdots\circ\mathcal{B}_{1}\circ\iota$ with equivariant linear maps $\iota:\mathcal{V}\to\mathcal{H}$ and $\pi:\mathcal{H}\to\mathcal{W}$.
Each block is equivariant: a rotation translates the argument of the synthesized fields, the point-wise product commutes with this translation, and the Haar projection~\eqref{eq:sm-wbb-analysis} intertwines it back to the coefficient action.

\begin{lemma}[Clebsch--Gordan trees span polynomial equivariants]
\label{lem:sm-cg-span}
Let $\mathcal{V},\mathcal{W}$ be finite direct sums of SO(3) irreps.
Every polynomial SO(3)-equivariant map $P:\mathcal{V}\to\mathcal{W}$ is a finite linear combination of \emph{tree maps}: compositions of equivariant linear leaf maps $\mathcal{V}\to V_{l}$, iterated binary couplings~\eqref{eq:sm-cg-map}, and a final equivariant linear map into $\mathcal{W}$, together with constant terms in the trivial isotypic part of $\mathcal{W}$.
\end{lemma}

\begin{proof}
Decomposing $P$ into homogeneous parts and comparing degrees in $P(tx)$ shows each part is equivariant; the degree-zero part is an invariant constant.
A homogeneous part of degree $d$ is, by polarization, the diagonal restriction of a symmetric $d$-linear equivariant map, hence of an element of $\mathrm{Hom}_{SO(3)}(\mathcal{V}^{\otimes d},\mathcal{W})$; it therefore suffices to span this space by trees, and, after decomposing $\mathcal{V}$ and $\mathcal{W}$ into irreps, to span $\mathrm{Hom}_{SO(3)}(V_{l_1}\otimes\cdots\otimes V_{l_d},V_L)$.
We induct on $d$.
For $d=1$ Schur's lemma leaves only multiples of the identity.
For $d>1$, decompose $X=V_{l_1}\otimes\cdots\otimes V_{l_{d-1}}$ into irreducible summands with equivariant projections $p_{J,r}:X\to V_J$ and inclusions $i_{J,r}$ (the equivariant maps extracting and identifying the $r$-th copy of $V_J$ in the decomposition of $X$ into irreducibles, with $\sum_{J,r}i_{J,r}p_{J,r}=\mathrm{id}_X$); any $T\in\mathrm{Hom}_{SO(3)}(X\otimes V_{l_d},V_L)$ writes as $T=\sum_{J,r}T\circ(i_{J,r}p_{J,r}\otimes\mathrm{id})$, and each factor $T\circ(i_{J,r}\otimes\mathrm{id})$ lies in $\mathrm{Hom}_{SO(3)}(V_J\otimes V_{l_d},V_L)$, which by the multiplicity-free Clebsch--Gordan decomposition is spanned by $C^{L}_{Jl_d}$ when the triple is admissible and vanishes otherwise.
The induction hypothesis expresses each $p_{J,r}$ by trees on the first $d-1$ factors; composing with $C^{L}_{Jl_d}$ yields trees on all $d$ factors.
Evaluating on the diagonal $(x,\ldots,x)$ recovers the homogeneous polynomial parts and proves the claim.
\end{proof}

\begin{theorem}[Completeness of the WBN]
\label{thm:sm-wbn-complete}
Let $\mathcal{V},\mathcal{W}$ be finite direct sums of SO(3) irreps and let $F:\mathcal{V}\to\mathcal{W}$ be continuous and SO(3)-equivariant.
For every compact $\Omega\subset\mathcal{V}$ and every $\varepsilon>0$ there exist a band limit $L_{\max}$, a width $C$, a depth $K$, and WBN parameters such that
\begin{equation}
    \sup_{x\in\Omega}
    \bigl\|\pi\circ\mathcal{B}_{K}\circ\cdots\circ\mathcal{B}_{1}\circ\iota(x)-F(x)\bigr\|
    <\varepsilon .
    \label{eq:sm-wbn-uat}
\end{equation}
Moreover, every \emph{polynomial} SO(3)-equivariant map is realized exactly, with zero error.
\end{theorem}

\begin{proof}
\emph{Step 1 (reduction to polynomial equivariants).}
Let $\widehat\Omega=SO(3)\,\Omega$, compact by compactness of the group.
By the Stone--Weierstrass theorem there is a polynomial map $R:\mathcal{V}\to\mathcal{W}$ with $\sup_{\widehat\Omega}\|R-F\|<\delta$.
Its Reynolds average $P(x)=\int_{SO(3)}g^{-1}R(gx)\,dg$ is polynomial and equivariant, and since $F(x)=\int g^{-1}F(gx)\,dg$,
\[
    \sup_{x\in\Omega}\|P(x)-F(x)\|
    \le \Bigl(\sup_{g}\|g^{-1}\|_{\mathcal{W}}\Bigr)\,\delta ,
\]
which is smaller than $\varepsilon$ for $\delta$ small.
It therefore suffices to realize polynomial equivariants exactly.

\emph{Step 2 (one block realizes one coupling).}
Suppose two hidden channels of the current feature hold covariants $T_1(x)\in V_{l_1}$ and $T_2(x)\in V_{l_2}$, and fix an admissible output degree $L$.
Since the coupling map~\eqref{eq:sm-cg-map} is nonzero, its coefficient array is not identically zero, so there exist frames $|k_1|\le l_1$, $|k_2|\le l_2$ with
\begin{equation}
    \langle l_1 k_1\, l_2 k_2|L\,\kappa\rangle\neq0,
    \qquad \kappa=k_1+k_2 .
    \label{eq:sm-frame-choice}
\end{equation}
Use the maps~\eqref{eq:sm-wbb-lin} to load the two channels into one product pair,
\[
    f^{\mathrm{L}}_{a}(g)=\sum_{m_1}(T_1)_{m_1}D^{(l_1)}_{m_1k_1}(g),
    \qquad
    f^{\mathrm{R}}_{b}(g)=\sum_{m_2}(T_2)_{m_2}D^{(l_2)}_{m_2k_2}(g),
\]
all other coefficients set to zero.
By the product identity~\eqref{eq:sm-wigner-product},
\[
    h_{ab}(g)
    =\sum_{L'}\langle l_1 k_1\, l_2 k_2|L'\kappa\rangle\,
    \bigl(C^{L'}_{l_1l_2}(T_1,T_2)\bigr)_{M}\,D^{(L')}_{M\kappa}(g),
\]
and the analysis~\eqref{eq:sm-wbb-analysis} with the orthogonality~\eqref{eq:sm-wigner-orth} extracts
\[
    z^{(L)}_{M,\kappa,ab}
    =\langle l_1 k_1\, l_2 k_2|L\,\kappa\rangle\,
    \bigl(C^{L}_{l_1l_2}(T_1,T_2)\bigr)_{M}.
\]
The nonzero scalar~\eqref{eq:sm-frame-choice} is absorbed into $B^{(L)}_{d,\kappa,ab}$ in~\eqref{eq:sm-wbb-out}; choosing $P^{(l)}$ to preserve the existing channels and one fresh output channel $d$ stores exactly the coupling $C^{L}_{l_1l_2}(T_1,T_2)$.

\emph{Step 3 (finite stacks realize all trees).}
Fix a polynomial equivariant $P$ and expand it by Lemma~\ref{lem:sm-cg-span} into finitely many tree maps.
Assign one hidden channel to every node of every tree; choose $L_{\max}$ no smaller than the largest degree appearing and $C$ large enough to hold all channels.
The embedding $\iota$ places the leaf maps.
Proceeding by induction on tree depth, each block writes all couplings of the next depth into fresh channels -- Step~2, applied with a separate product pair $(a,b)$ and output channel per node, the mixing coefficients $B$ set to ignore all other pairs -- while preserving the channels already filled.
After finitely many blocks all nodes are present; the readout $\pi$ forms the finite linear combination of tree outputs, and the bias $q$ supplies the constant terms.
The stack realizes $P$ exactly, and with Step~1 the theorem follows.
\end{proof}

\begin{remark}[Frame truncation and sphere projection]
\label{rem:sm-wbn-truncation}
The full frame range enters the proof only through the choice~\eqref{eq:sm-frame-choice}.
Truncating the frames to $|k|\le k_{\max}$ restricts that choice, and at $k_{\max}=0$ it forces $k_1=k_2=\kappa=0$: the surviving coefficient $\langle l_1 0\, l_2 0|L 0\rangle$ vanishes whenever $l_1+l_2+L$ is odd, so a sphere-restricted bilinear nonlinearity cannot realize any parity-odd coupling, however deep the stack -- the incompleteness discussed in main-text Section~\ref{sec:s2-method}.
Already at $k_{\max}=1$ nonzero parity-odd choices become available, e.g.\ $\langle 1\,1\,1\,{-1}|1\,0\rangle=1/\sqrt{2}$ for the cross product $1\otimes1\to1$; the hypothesis of Theorem~\ref{thm:sm-wbn-complete}, however, requires the full range.
\end{remark}

\FloatBarrier
\section{Training and systems methods}
\label{sec:sm-training-systems}

\subsection{Training objective}
\label{sec:sm-training-objective}

DPA4 is trained as a conservative interatomic potential.
The neural network predicts a scalar total energy, and forces are obtained by differentiating this energy with respect to atomic positions.
The reported training runs use MAE losses with vector-norm force residuals.
For a mini-batch of configurations
$b=1,\ldots,B$, with $N_b$ atoms in configuration $b$, the objective is
\begin{equation}
    \mathcal{L}
    =
    \lambda_E
    \frac{1}{B}\sum_{b=1}^{B}
    \frac{|E_{\Theta,b}-E_b|}{N_b}
    +
    \lambda_F
    \frac{1}{\sum_{b=1}^{B}
        N_b} \sum_{b=1}^{B}\sum_{i=1}^{N_b} \left\| \mathbf{F}_{\Theta,bi}-\mathbf{F}_{bi} \right\|_2 + \lambda_{\Pi} \frac{1}{B}\sum_{b=1}^{B} \frac{\left\|\Pi_{\Theta,b}-\Pi_b\right\|_1}{9N_b}.
    \label{eq:sm-loss}
\end{equation}
Here $E_{\Theta,b}$, $\mathbf{F}_{\Theta,bi}$ and $\Pi_{\Theta,b}$ denote DPA4 predictions, while $E_b$, $\mathbf{F}_{bi}$ and $\Pi_b$ denote reference DFT labels.
The force residual is treated as a three-dimensional vector for each atom: the Euclidean norm is taken before averaging over atoms.
Energy and virial residuals use the MAE form with per-atom normalization.
The benchmark-specific weights, batch sizes and training lengths are listed in Tables~\ref{tab:sm_main_ablation_config}--\ref{tab:sm_matbench_config}.

\subsection{Warmup--stable--decay learning-rate schedule}
\label{sec:sm-wsd-schedule}

For long training runs, DPA4 uses a warmup--stable--decay schedule in which the learning rate first increases linearly, remains constant for the main training phase and is annealed only near the end of the run~\cite{wen2024wsd}.
Let $T$ be the total number of optimization steps, $T_{\mathrm{w}}$ the warmup length, $T_{\mathrm{d}}$ the decay length and $T_{\mathrm{s}}=T-T_{\mathrm{w}}-T_{\mathrm{d}}$ the stable length.
The
schedule used in the reported DPA4 training runs is
\begin{equation}
    \alpha(t)=
    \begin{cases}
        \alpha_{\mathrm{w}}
        +(\alpha_0-\alpha_{\mathrm{w}})t/T_{\mathrm{w}},
         & 0\le t<T_{\mathrm{w}},                             \\[3pt]
        \alpha_0,
         & T_{\mathrm{w}}\le t<T_{\mathrm{w}}+T_{\mathrm{s}}, \\[3pt]
        \alpha_{\mathrm{d}}(\tau),
         & T_{\mathrm{w}}+T_{\mathrm{s}}\le t<T,              \\[3pt]
        \alpha_{\min},
         & t\ge T,
    \end{cases}
    \label{eq:sm-wsd-piecewise}
\end{equation}
where $\alpha_{\mathrm{w}}$ is the initial warmup learning rate,
$\alpha_0$ is the stable-phase learning rate,
$\alpha_{\min}$ is the final learning rate and
\begin{equation}
    \tau
    =
    \operatorname{clip}
    \left(
    \frac{t-T_{\mathrm{w}}-T_{\mathrm{s}}}{T_{\mathrm{d}}},
    0,1
    \right).
    \label{eq:sm-wsd-tau}
\end{equation}
with cosine annealing in the decay phase,
\begin{equation}
    \alpha_{\mathrm{d}}(\tau)
    =
    \alpha_{\min}
    +
    \frac{\alpha_0-\alpha_{\min}}{2}
    \left(1+\cos \pi\tau\right).
    \label{eq:sm-wsd-cosine}
\end{equation}
Thus the stable phase carries most optimization steps, whereas the final cosine decay suppresses high-learning-rate oscillations before checkpoint selection.

\subsection{HybridMuon optimizer}
\label{sec:sm-hybrid-muon}

DPA4 is optimized with HybridMuon, a matrix-aware hybrid optimizer adapted from Muon~\cite{jordan2024muon} and scalable-Muon training studies~\cite{liu2025muon}.
The design separates two classes of trainable parameters.
Matrix-valued hidden transformations are routed to Muon, whereas biases, normalization scales, one-dimensional parameters and explicitly marked auxiliary parameters are routed to Adam or decoupled-weight-decay Adam.
This static routing is built on the first optimizer step and remains fixed during training, so the update rule for each parameter is independent of the current gradient value.

\paragraph{Matrix views and slice mode.}
Let a trainable tensor have an effective shape obtained by removing singleton dimensions.
HybridMuon interprets a rank-two effective shape as one matrix.
For higher-rank equivariant tensors, DPA4 uses slice mode: the leading dimensions index independent blocks, and Muon is applied separately to every trailing $(m,n)$ matrix.
Thus, for an effective shape $(b_1,\ldots,b_q,m,n)$, the
optimizer constructs
\begin{equation}
    B=\prod_{a=1}^{q}b_a \quad\text{independent matrix blocks}\quad G_t^{(b)}\in\mathbb{R}^{m\times n}, \qquad b=1,\ldots,B.
    \label{eq:sm-muon-slice-view}
\end{equation}
This choice is important for SO(2)-structured equivariant weights: degree and order-indexed blocks are updated independently, rather than being flattened into one large matrix that would mix unrelated representation strata.
In the degree-wise SO(3) channel maps used by the SO(2) pre- and post-projections and by the equivariant FFN, the weight tensor has shape $(L+1,C_{\mathrm{in}},F C_{\mathrm{out}})$; slice mode therefore applies Muon separately to the $(C_{\mathrm{in}},F C_{\mathrm{out}})$ matrix of each degree $l$.
Local SO(2) linear maps keep separate matrix parameters for the $m=0$ block and the constrained $|m|>0$ blocks, so the optimizer acts on these structured matrix blocks without collapsing distinct equivariant subspaces.

\paragraph{Muon update.}
For each Muon-routed block, the optimizer maintains a momentum buffer and forms
a Nesterov-style update
\begin{align}
    M_t^{(b)} & = \beta M_{t-1}^{(b)}+(1-\beta)G_t^{(b)}, \\ U_t^{(b)} & = \beta M_t^{(b)}+(1-\beta)G_t^{(b)}.
    \label{eq:sm-muon-nesterov}
\end{align}
The matrix $U_t^{(b)}$ is then orthogonalized by a Newton--Schulz polar iteration.
For the standard square or batched path, the iteration starts from
$X_0=U_t^{(b)}/\|U_t^{(b)}\|_{\mathrm{F}}$ and applies
\begin{equation}
    X_{k+1} = a_k X_k + \left(b_k A_k+c_k A_k^2\right)X_k, \qquad A_k=X_kX_k^{\mathsf{T}}.
    \label{eq:sm-newton-schulz}
\end{equation}
DPA4 uses a two-stage schedule: eight fast iterations with $(a,b,c)=(3.4445,-4.7750,2.0315)$ followed by two Newton polishing iterations with $(a,b,c)=(2,-1.5,0.5)$.
The resulting polar factor is denoted $Q_t^{(b)}$.

\paragraph{Rectangular Gram path.}
For rectangular matrices, HybridMuon uses a compiled Gram Newton--Schulz path following the fast polar-decomposition formulation used for Muon~\cite{dao2026gramnewtonschulz}.
The matrix is oriented so that $m\le n$, normalized in single precision and iterated in half precision using a fixed Polar-Express coefficient schedule.
Since the iteration depends on $XX^{\mathsf{T}}\in\mathbb{R}^{m\times m}$, rectangular blocks with the same smaller dimension can be column-padded, concatenated and orthogonalized in a single grouped call.
Padding only the larger dimension preserves the Frobenius
norm and the Gram matrix,
\begin{equation}
    [X\;0][X\;0]^{\mathsf{T}}=XX^{\mathsf{T}},
    \label{eq:sm-gram-padding}
\end{equation}
so truncating the padded columns after the iteration exactly recovers the
unpadded result while reducing the number of small GPU launches.

\paragraph{Update-RMS matching and Adam-family path.}
In the default match-RMS mode, the Muon update for an $m\times n$ block is
scaled as
\begin{equation}
    \Delta W_t^{(b)}
    =
    -\alpha_t\,\gamma\,\sqrt{\max(m,n)}\,Q_t^{(b)},
    \qquad \gamma=0.18.
    \label{eq:sm-match-rms}
\end{equation}
This coefficient follows the update-RMS calibration used in scalable Muon training~\cite{liu2025muon}: it brings the per-element magnitude of the orthogonalized matrix update onto the same learning-rate scale as AdamW-like updates.
Parameters routed to Adam use first and second moments in single precision, with bias correction and $\epsilon=10^{-20}$.
Decoupled weight decay is applied to Muon-routed matrices and to decay-enabled Adam-routed tensors, but not to one-dimensional Adam-routed parameters.

\subsection{Magma-lite update damping}
\label{sec:sm-magma-lite}

Conservative energy-gradient training produces gradients that include mixed coordinate--parameter derivatives and can exhibit large block-to-block variation.
DPA4 therefore augments the Muon path with a deterministic Magma-lite damping rule, adapted from momentum-aligned gradient masking~\cite{joo2026magma}.
For each Muon block, let $G_t^{(b)}$ be the current gradient and $M_t^{(b)}$ the momentum buffer after the update in Eq.~\eqref{eq:sm-muon-nesterov}.
The block alignment is
\begin{equation}
    \chi_t^{(b)}
    =
    \frac{
        \left\langle M_t^{(b)},G_t^{(b)}\right\rangle_{\mathrm{F}}
    }{
        \|M_t^{(b)}\|_{\mathrm{F}}\,
        \|G_t^{(b)}\|_{\mathrm{F}}+\varepsilon
    },
    \qquad
    \chi_t^{(b)}\in[-1,1].
    \label{eq:sm-magma-cosine}
\end{equation}
The score is mapped through a temperature-scaled sigmoid and stretched to
$[0,1]$,
\begin{equation}
    r_t^{(b)}
    =
    \operatorname{clip}
    \left[
        \frac{
            \sigma(\chi_t^{(b)}/\tau)-\sigma(-1/\tau)
        }{
            \sigma(1/\tau)-\sigma(-1/\tau)
        },
        0,1
        \right],
    \qquad \tau=2.
    \label{eq:sm-magma-score}
\end{equation}
An exponential moving average gives
\begin{equation}
    u_t^{(b)}
    =
    \rho_{\mathrm{ema}} u_{t-1}^{(b)}
    +(1-\rho_{\mathrm{ema}})r_t^{(b)},
    \qquad \rho_{\mathrm{ema}}=0.9,
    \label{eq:sm-magma-ema}
\end{equation}
and the final Muon update is rescaled by
\begin{equation}
    \Delta W_t^{(b)}
    \leftarrow
    \left[s_{\min}+(1-s_{\min})u_t^{(b)}\right]\Delta W_t^{(b)},
    \qquad s_{\min}=0.1.
    \label{eq:sm-magma-lite}
\end{equation}
This rule differs from stochastic masking: all blocks remain active, and poorly aligned updates are continuously damped rather than randomly skipped.
The nonzero lower bound prevents the optimizer from freezing a block entirely, which is important for MLIP training where force labels can make the alignment temporarily noisy without implying that the corresponding representation block should stop learning.

\subsection{Compiled conservative energy-gradient training}
\label{sec:sm-compiled-training}

The conservative energy-gradient path is compiled separately from direct-force or density-denoising modes.
Force matching requires differentiating through
\begin{equation}
    \mathbf{F}_{\Theta}
    =
    -\frac{\partial E_{\Theta}}{\partial R},
    \qquad
    \frac{\partial\mathcal{L}}{\partial\Theta}
    \supset
    \frac{\partial^2E_{\Theta}}{\partial R\,\partial\Theta}.
    \label{eq:sm-double-backward}
\end{equation}
A standard compiled forward/backward stack captures the forward graph and its first reverse-mode derivative, but it does not directly expose a nested coordinate derivative inside the compiled region.
DPA4 resolves this by first executing the energy-to-force derivative during symbolic tracing and then lowering the resulting tensor graph with PyTorch Inductor~\cite{ansel2024pytorch2}.

\paragraph{Tracing the conservative lower graph.}
The compiled function wraps the lower energy computation as a tensor-only map: from extended coordinates, atom types, neighbor indices and optional conditioning tensors to energies, forces and virials.
Before entering the traced function, the extended coordinates are rebound to a fresh leaf tensor.
This restart removes any upstream graph carried by data loading or neighbor construction while preserving a well-defined coordinate endpoint for $\partial E/\partial R$.
During training, the force construction keeps the coordinate-derivative graph alive; the outer loss backward pass can therefore differentiate the force residual into the model parameters, realizing the mixed derivative in Eq.~\eqref{eq:sm-double-backward}.

\paragraph{Symbolic tracing and higher-order differentiability.}
The implementation traces the conservative lower graph with symbolic shapes and real tensor inputs.
Real inputs are needed because compact edge construction contains data-dependent operations; after that control flow is resolved, the runtime dimensions are represented symbolically.
The trace uses a small five-frame representative batch chosen to avoid known symbolic-dimension collisions with singleton axes, charge--spin width, Cartesian coordinates and virial components.
The SiLU backward operation is decomposed into elementary pointwise operations before tracing, so the compiler receives an explicit first-derivative graph when the optimizer later requests a second derivative.

\paragraph{Preserving the force-loss gradient.}
When the coordinate derivative is traced in training mode, autograd inserts detach nodes around saved forward activations.
In the traced FX graph these nodes would become ordinary tensor operations and would sever the path from the force loss back to the parameters.
DPA4 removes only the detach nodes matching the saved-tensor topology and keeps user-intended detach operations unchanged.
The edited graph is then rebuilt into a fresh FX graph before compilation, so all compiler passes see a consistent node topology.

\paragraph{Dynamic edge representation and cache structure.}
The padded DeePMD neighbor list is converted inside the graph into a compact edge list.
Edge vectors are formed by differentiable indexed selection from the extended coordinate tensor, and a single masked sentinel edge is appended to every batch.
The sentinel edge guarantees a nonempty edge tensor under symbolic shapes while contributing exactly zero to downstream reductions.
Compiled callables are cached by graph topology, including training versus evaluation, the presence of atomic virial outputs and coordinate-correction inputs.
This multi-slot cache avoids recompilation when training is periodically interrupted by validation, while keeping distinct output signatures in separate compiled graphs.

\paragraph{Inductor configuration.}
The traced graph is lowered with dynamic-shape compilation.
The compiled path uses deterministic compilation settings rather than autotuning, enables shape padding for fluctuating symbolic dimensions and disables compiler features that interfere with higher-order autograd metadata or produce unstable large fused reduction kernels for this higher-order graph.
Training and evaluation use separate fusion limits because the training graph contains the second-derivative branch whereas the evaluation graph does not.
This systems design preserves the scalar energy-to-force relation while allowing the conservative training path to remain inside compiled GPU code.

